\section{Conclusion}
\label{sec:geometry-first-conclusion}

This work develops a geometry-first workflow in which the physical degrees of
freedom are outputs of constrained reduction rather than inputs chosen as
fundamental fields.  BHSM supplies the organizing principle: vary first,
enforce constraints and boundary data, quotient gauge redundancy, construct
the physical Hessian, and only then identify the response variables that
propagate to observables.

Applied to closed hyperspherical cosmology, that procedure yields a physical
$n=2$ reduced state, an induced metric response, a gauge-invariant
luminosity-distance kernel, and a Jacobi/Sachs optical transport system for
redshift, convergence, and shear.  Unresolved topographic structure can be
propagated statistically, leading to a low-rank cross-observable covariance
prediction when only a small number of physical texture channels are active.

Effective-field descriptions remain valuable comparison tools, but they are
not taken here as the defining ontology of the construction.  The
curvature-dependent compatibility failure encountered in one effective
representation and the rejection of several preregistered foreground models
illustrate the practical importance of keeping representation distinct from
physical reduction.

The remaining challenge is upstream rather than observational: derive the
common-action normalization and physical source statistics of the local
topographic response.  The downstream chain from physical response to photon
transport and observable kernels is now sufficiently explicit to make those
future source calculations empirically consequential.

The central proposal is therefore methodological and falsifiable: geometry and
physical reduction should precede field representation, and the surviving
reduced structure must earn its status by predicting several observables with
one common response architecture.
